\section{Uniform collision clouds and confidence geometry}\label{sec:clouds}
Put $\tau=\min(1,d_0/8)$, and write $Q_H=z+[-H,H]^2$ for a working
square. For a capped distance radius $c>0$, launch in
$Q_{H+c+2D_0+2\tau}$ and keep first impacts up to time $\tau$.
The square is translated with the observed flight; its area, denoted $S$,
is known. Only first impacts in $Q_{H+c+D_0+\tau}$ are retained.
Every body meeting $Q_{H+c}$ and its outward collar of width $\tau/2$
are contained in the launch square, and its entire boundary is retained.

\begin{lemma}[First-hit minorization]\label{lem:flux-minor}
For every boundary arc $I$ of such a body and every launch attempt,
\begin{equation}\label{eq:minor}
 \Prob\{\hbox{the first impact before }\tau\hbox{ lies in }I\}
                         \ge \frac{\tau|I|}{8\pi S}.
\end{equation}
The bound does not condition on a free initial position or on a collision.
\end{lemma}
\begin{proof}
At $x\in I$ let an incoming direction make angle $\phi\in[-\pi/3,\pi/3]$
with the inward normal, and put $q=x-rv$, $r\in[\tau/4,\tau/2]$.
The segment from $q$ to $x$ lies outside the source body by its supporting
half-plane. It cannot meet another body since its length is below $d_0$.
Thus $x$ is its first impact. The coordinates $(s,\phi,r)$ give phase
volume $\cos\phi\,\dd s\dd\phi\dd r$; the first-hit parametrization is
injective for these states. Dividing by $2\pi S$ gives
$\sqrt3\tau|I|/(8\pi S)$, which implies \eqref{eq:minor}.
Solid launch failures are included in the normalization $S$.
\end{proof}

Only finitely many bodies meet this working region. For example, put
$H'=H+c+3D_0+2\tau+d_0$ and
\[
 M_0=\left\lceil64(1+H'/d_0)^2\right\rceil.
\]
Choose one point from each relevant body; their mutual distances are at
least $d_0$. Disjoint disks of radius $d_0/2$ around these points fit in
the square enlarged by $D_0+d_0$, proving the conservative bound $M_0$.
The circumference of a convex body of diameter $D_0$ is at most $\pi D_0$.

\begin{lemma}[Coverage, clustering and Hausdorff enclosures]\label{lem:cloud}
Choose
\[
 0<q<\min\{d_0/40,1,\pi/\kappa_+\},\qquad
 \varepsilon_x<\min\{d_0/100,q/10\}.
\]
If
\begin{equation}\label{eq:sample-budget}
 n\ge\frac{16\pi S}{\tau q}
       \log\frac{2M_0(1+\pi D_0/q)}{\zeta},
\end{equation}
then, with probability at least $1-\zeta$, the following hold simultaneously.
Every relevant boundary has a retained impact within arclength distance
$q$ of each point. Connect reported positions at Euclidean distance below
$d_0/3$. No component joins different physical bodies, and the cloud of
each fully observed body is connected. Its convex hull $P$ satisfies
\begin{equation}\label{eq:hull}
                      d_H(P,C)\le e_q:=\kappa_+q^2/2+\varepsilon_x.
\end{equation}
Seeds at the encountered collisions may be included with the same
localization tolerance.
\end{lemma}
\begin{proof}
Partition each relevant boundary into arcs with lengths between $q/2$
and $q$; this is possible since its circumference is at least
$2\pi/\kappa_+>2q$. There are at most $1+\pi D_0/q$ arcs per body.
By \eqref{eq:minor}, the probability that any one arc is empty after
$n$ independent attempts is at most $\exp(-n\tau q/(16\pi S))$.
A union bound gives the assertion about coverage. Every point is within
arclength $q$ of a sample in its partition arc; consecutive retained
samples have arclength gaps at most $2q$. Localization therefore still
connects consecutive samples at distance below $d_0/3$. Points from
different bodies remain more than $d_0-2\varepsilon_x>d_0/3$ apart.
Partial clouds from other bodies cannot merge into a complete cloud.

At a support point $x(0)$ in direction $n$, use unit-speed arclength and
$|x''|\le\kappa_+$. For $|s|\le q$,
\[
 0\le n\cdot(x(0)-x(s))\le\kappa_+s^2/2,
\]
since $n\cdot x'(0)=0$. The support function of the exact sample hull
therefore lies between that of $C$ and that function minus
$\kappa_+q^2/2$. Localization changes a hull's support function by at
most $\varepsilon_x$. The support characterization of Hausdorff distance
proves \eqref{eq:hull}, uniformly in direction.
\end{proof}

This argument is a finite-sample form of the boundary random-polytope
mechanism; see \cite{Schneider,PSSW}. What is needed here, and not supplied
by a boundary-sampling oracle, is the physical first-hit minorization,
the treatment of unknown neighboring bodies, and the conditional use of
one finite cloud after an adaptively located flight.

\subsection{No unobserved-solid query}\label{sec:distance}
Take the union $P_{\rm all}$ of convex hulls of all retained components,
including partial exterior components. Let
$d_P(x)=\min\{c,\dist(x,P_{\rm all})\}$, with value $c$ for the empty
union. On the event of Lemma~\ref{lem:cloud}, for every $x\in Q_H$,
\begin{equation}\label{eq:distance-cloud}
 \left|d_P(x)-\min\{c,\dist(x,\mathcal O)\}\right|\le e_q.
\end{equation}
Indeed a hull formed from one body's reported points lies in its
$\varepsilon_x$-neighborhood, so it cannot create a spurious solid farther
than that error from the true solid. If the true distance is below $c$,
a nearest body meets $Q_{H+c}$, is completely covered, and has the
Hausdorff enclosure \eqref{eq:hull}. If the distance is at least $c$,
only the preceding one-sided bound is needed after capping.
Thus \eqref{eq:distance-cloud} controls all query positions at once.
There is no additional pointwise union bound and no distance oracle.

\subsection{Two derivatives and intrinsic coordinates}
Translate a complete cloud by one observed seed. Its true support
function has a $C^3$ bound $B_3=2+M_6+2D_0$, enlarged if necessary to
absorb the seed error. Fix the nonnegative kernel
\[
 \varphi(t)=\frac{315}{256}(1-t^2)^4\one_{[-1,1]}(t),\qquad
 \int\varphi=1,
\]
and its periodic rescaling $\varphi_h$. Set
$\widehat p=p_P*\varphi_h$. This is the support function of an average
of rotations of $P$, and hence remains a convex support function. Let
$C_\varphi=1+\|\varphi'\|_1+\|\varphi''\|_1$; the numerical bound
$C_\varphi\le512$ suffices below. For $h\le1$,
\begin{equation}\label{eq:mollify}
 \|\widehat p-p\|_{C^2}
       \le C_\varphi e_qh^{-2}+B_3h.
\end{equation}
For the first term differentiate the kernel up to twice. For the second,
write $p^{(j)}(\theta-hu)-p^{(j)}(\theta)$ as an integral of
$p^{(j+1)}$, $j\le2$, and use $\int|u|\varphi(u)\dd u\le1$.
This proves the displayed inequality without differentiating a noisy hull.

For $0<\nu<1$ choose
\begin{equation}\label{eq:explicit-cloud-choice}
 \begin{gathered}
 h=\frac{\nu}{8B_3},\qquad
 E_\nu=\frac{\nu^3}{4096C_\varphi B_3^2},\\
 q=\min\{d_0/80,1/2,\pi/(2\kappa_+),
                         \sqrt{E_\nu/(4\kappa_+)}\},\\
 \varepsilon_x\le\min\{E_\nu/8,q/20,d_0/200\}.
 \end{gathered}
\end{equation}
For the small-$\nu$ range where the last entry fixes $q$,
$e_q\le E_\nu/4$. Equations \eqref{eq:mollify} and
\eqref{eq:explicit-cloud-choice} give an error smaller than $\nu/4$.
The finite remaining range is handled by the smaller fixed $q$.
The sample count \eqref{eq:sample-budget} is
$C\nu^{-3/2}\log(C/(\nu\zeta))$, with all constants determined by
bounded local size and the numerical priors. Localization tolerance is
\emph{at most} the displayed value; smaller tolerated errors cost more,
not less, measurement precision.

For small enough $\nu$, $\widehat p+\widehat p''$ stays above
$1/(2\kappa_+)$. The recovered curve
\[
 \widehat x(\theta)=\widehat p(\theta)n(\theta)
                         +\widehat p'(\theta)t(\theta)
\]
is regular and strictly convex. Its radius of curvature differs by
$O(\nu)$ from the true one. Since
$\dd s=(p+p'')\dd\theta$, integration gives intrinsic arclength errors
$O(\nu)$ on bounded arcs. Inverting this positive speed and using the
true $C^3$ bound gives $O(\nu)$ errors for curvature in the common
intrinsic parameter. The Frenet equations then give intrinsic $C^2$
arc error $O(\nu)$. Coordinate rotation to a graph whose tangent has
fixed slack gives $C^2$ graph enclosures with the same order. No
arclength measurement was requested from the apparatus.

All estimates are uniform conditional on the preceding history: after
conditioning, the local square and unknown relevant bodies are fixed,
and only fresh auxiliary launches are used. Quadrature for the convolution
and arc integrals is over explicit hull support functions and a fixed
kernel. Its deterministic errors can be included in the budget
$\nu/4$, since their moduli and the finite hull are known. This proves
the geometric part of Theorem~\ref{thm:probe-main}.
